% Part: sets-functions-relations
% Chapter: arithmetization
% Section: rationals

\documentclass[../../../include/open-logic-section]{subfiles}

\begin{document}

\olfileid{sfr}{arith}{rat}
\olsection{Saka $\Int$ menyang $\Rat$}

Kita lagi wae weruh carane mbangun wilangan wutuh saka wilangan
natural kanthi nggunakake sawatara teori himpunan naif. Saiki kita
bakal weruh carane mbangun wilangan rasional saka wilangan wutuh
kanthi cara kang meh padha. Pangamatan wiwitan yaiku:
\begin{enumerate}
	\item Saben wilangan rasional bisa ditulis kanthi wujud $\nicefrac{i}{j}$,
	kanthi $i$ lan $j$ loro-lorone wilangan wutuh nanging $j$ ora nol.
	\item Informasi kang dikodeake ing ekspresi $\nicefrac{i}{j}$
	uga bisa dikodeake kanthi pasangan urut $\tuple{ i, j}$.
\end{enumerate}
Cara kang cetha yaiku nganggep wilangan rasional \emph{minangka}
pasangan urut kang dijupuk saka $\Int \times (\Int \setminus \{0_\Int\})$.
Nanging kaya sadurunge, cara iki rada naif banget, amarga kita kepengin
$\nicefrac{3}{2} = \nicefrac{6}{4}$, kamangka $\tuple{ 3, 2}\neq \tuple{ 6,
4}$. Luwih umum, kita kepengin sarat iki:
\[
	\nicefrac{a}{b} = \nicefrac{c}{d} \text{ yen lan mung yen } a \times d = b \times c
\]
Kanggo ngolehake iki, kita netepake {relasi ekuivalensi} ing $\Int \times
(\Int \setminus \{0_\Int\})$ kaya mangkene:
\[
	\tuple{ a, b }\Ratequiv \tuple{ c, d} \text{ yen lan mung yen }a \times d = b \times c
\]
Kita kudu mriksa yen iki relasi ekuivalensi. Iki meh padha banget karo
kasus $\Intequiv$, mula kita pasrahake minangka gladhen.
\begin{prob}
Tuduhna yen $\Ratequiv$ iku relasi ekuivalensi.
\end{prob}
Nanging relasi iki marakake kita bisa menehi definisi iki:
\begin{defn}
Wilangan rasional yaiku kelas-kelas ekuivalensi miturut $\Ratequiv$ saka
pasangan wilangan wutuh (kang unsur kapindhone ora nol). Tegese, $\Rat =
\equivclass{(\Int \times (\Int\setminus \{0_\Int\}))}{\Ratequiv}$.
\end{defn}

Kaya tumrap wilangan wutuh, kita uga kepengin netepake sawetara operasi
dhasar. Yen $\equivrep{i,j}{\Ratequiv}$ iku kelas ekuivalensi miturut
$\Ratequiv$ kang duwe $\tuple{i, j}$ minangka !!a{element}, kita netepake:
\begin{align*}
	\equivrep{a, b}{\Ratequiv} + \equivrep{c, d}{\Ratequiv} &= \equivrep{ad + bc,  bd}{\Ratequiv}\\
	\equivrep{a, b}{\Ratequiv} \times \equivrep{c, d}{\Ratequiv} &= \equivrep{a  c, b d}{\Ratequiv}.
\intertext{Kanggo netepake $r \leq s$ ing wilangan rasional iki, kita nggunakake
kasunyatan yen $r \le s$ yen lan mung yen $s - r$ ora negatif, yaiku $s - r$
bisa ditulis minangka $\nicefrac{i}{j}$ kanthi $i$ ora negatif lan $j$~positif:}
	\equivrep{a, b}{\Ratequiv} \leq \equivrep{c, d}{\Ratequiv} &\text{ yen lan mung yen }
	\equivrep{c, d}{\Ratequiv} - \equivrep{a, b}{\Ratequiv} =
	\equivrep{i_\Int, j_\Int}{\Ratequiv}
\end{align*}
kanggo sawatara $i \in \Nat$ lan $0 \neq j \in \Nat$.
Cathetan koreksi sumber OLP-0043: ing ukara sadurunge, sumber Inggris
sapisan mbalikke urutan suku ing bedane; terjemahan iki manut urutan
kang bener ing rumus.

Banjur kita kudu mriksa yen definisi-definisi iki tumindak kaya kang
\emph{kudune}; lan rinciane dipindhah menyang \olref[check]{sec}. Nanging
pancen mangkono!{} Pungkasan, kita kepengin cara kanggo nganggep wilangan
wutuh \emph{minangka} wilangan rasional; mula kanggo saben $i \in \Int$,
kita netepake $i_\Rat = \equivrep{i, 1_\Int}{\Ratequiv}$. Maneh, kita
mriksa yen kabeh iki tumindak kanthi bener ing \olref[check]{sec}.

\begin{prob}
Tuduhna yen $(i + j)_\Rat = i_\Rat+ j_\Rat$ lan $(i \times j)_\Rat =
i_\Rat \times j_\Rat$ lan $i \leq j \liff i_\Rat \leq j_\Rat$, kanggo saben
$i, j \in \Int$.
\end{prob}

\end{document}
